%\newcommand{\diff}[1]{\dot{#1}}
\newcommand{\soln}{\gamma}
\newcommand{\solnt}{g}
\newcommand{\dsoln}{\solnt'}
\newcommand{\dsolnt}{\solnt'}
\newcommand{\soltest}{\phi}
\newcommand{\tsoltest}{{{\tilde{\phi}}}}

%The target dimension
\newcommand{\tdim}{{n}}

%Euer \g gleichsetzen mit  \soln...
\renewcommand{\g}{{\soln}}
\newcommand{\gn}{{\gamma_0}}

\section{Introduction}
To find nice geometric representatives within a given knot class, several new energies have been invented 
in the last two decades. 
The earliest knot energy for smooth curves  was the so-called \emph{M\"obius energy} introduced 
by Jun O'Hara
in \cite{OH91}, %\cite{OH92},\cite{OH94}
\[
 E^{(2)}(\soln) := \intl_{\R/\Z}\intl_{-1/{2}}^{1/{2}} 
  \brac{\fracm{\abs{\soln(u+w)-\soln(u)}^2}-\fracm{d_{\soln}(u+w,u)^2}}\ \abs{\soln'(u+w)}\abs{\soln'(u)}\ 
 \d w\d u
\]
which was later on extended to the family of energies
% \[
%  E^{(\alpha)}(\soln) := \intl_{\R/\Z}\intl_{-\fracm{2}}^{\fracm{2}} 
%   \brac{\fracm{\abs{\soln(u+w)-\soln(u)}^\alpha}-\fracm{D_{\soln}(u+w,u)}^\alpha}\ \abs{\soln'(u+w)}\abs{\soln'(u)}\ 
%  dw\ du
% \]
% for $\alpha \in [2,3)$ in \cite{OHara1992a} and even
\[
E^{(\alpha,p)}(\soln) := \intl_{\R/\Z}\intl_{-1/{2}}^{1/{2}} 
  \brac{\fracm{\abs{\soln(u+w)-\soln(u)}^\alpha}-\fracm{d_{\soln}(u+w,u)^\alpha}}^p\ \abs{\soln'(u+w)}\abs{\soln'(u)}\ 
 \d w\d u
\]
for $\alpha, p \in [1,\infty)$, see~\cite{OHara1994}.
Here $d_\g(u+w,u)$ denotes the intrinsic distance between $\g(u+w)$ and $\g(u)$ on the curve~$\g$.
More precisely,
\begin{equation}\label{eq:intrinsic}
 d_\g(u+w,u):= \min\big(\length(\g|_{[u,u+w]}),\length(\g)-\length(\g|_{[u,u+w]})\big)
\end{equation}
provided $\abs w\le\tfrac12$ where
$%\label{eq:length}
 \length(\g) := \int_0^1 \abs{\dg(\theta)}\d\theta
$
is the length of~$\g$.


Crucially using the M\"obius invariance of this knot energy, Michael Freedman, Zheng-Xu He, and Zhenghan Wang \cite{Freedman1994} were able to show that there are minimizers of the M\"obius energy within every 
prime knot class and that these are in fact of class $C^{1,1}$. More precisely, they
could show that if $\g$ is a local minimizer with respect to the $L^\infty$-topology, and 
if $\g$ is parametrized by arc length, then $\g$ is $C^{1,1}$.
Together with a bootstrapping argument due to He \cite{He2000}, one then obtains that local 
minimizers of the M\"obius energy are smooth, also see~\cite{Reiter2009a}.

Unfortunately, motivated by numerical evidence, Rob Kusner and John Sullivan were led to conjecture 
that there are no minimizers within composite knot classes \cite{Kusner1997}. In contrast to that, there
are minimizers of the energies $E^{(\alpha,p)}$ in the case that $jp>2$ as shown in \cite{OHara1994}. 

In this article, we will prove that even only stationary points of the M\"obius energy are of class
$C^\infty$ under the mildest condition one can think of: that $E^{(2)}(\g)$
is finite - an assumption which, as shown in a recent work of the first author \cite{Blatt2010a}, is equivalent to 
assuming that $\g$ is an injective curve of class $H^{\frac 3 2}(\R / \Z, \R ^n)$.
Our motivation to do so is twofold:
First of all, of course, 
this is a much stronger result than the smoothness of 
local minimizers as stated above.
Secondly, the M\"obius invariance is
not essential for proving smoothness of local minimizers as we do not need it in our arguments.
Thus there is the chance
to study other, possibly not M\"obius invariant, critical knot energies, using the techniques
developed in this article.
Additionally, our arguments are not restricted to the
one-dimensional situation but can be applied to arbitrary dimensions.
 
The price we pay is that, instead of the very appealing geometric argument in \cite{Freedman1994}, we have to adapt some 
sophisticated techniques originally developed  by Tristan Rivi\`ere and Francesca Da Lio \cite{DR11Sphere,DR11Man,DL11MannD}
and the third author \cite{SchikorranharmSphere,SchikorraIntBoundFrac11} to deal with $\frac n2$- harmonic maps into manifolds.

The first task in order to prove this result, is to derive the Euler-Lagrange equation for such stationary
points. In \cite{Freedman1994}, it was shown that for simple closed curves 
$\g \in C^{1,1}(\mathbb R / \mathbb Z, \mathbb R^n)$ and $h\in C^{1,1}(\mathbb R / \mathbb Z, \mathbb R^n)$
we have
\begin{align*}
 \delta E ^{(2)}(\g;h) &:= \lim_{\tau \searrow 0 } \frac {E^{(2)}(\g+\tau h) - E^{(2)}(\g)}\tau \\
 &= \int_{\mathbb R / \mathbb Z} \int_{\mathbb R / \mathbb Z} 
  \left(\frac {\g'(u) h'(u) - \frac {\left\langle \g(u)- \g(v), h(u)-h(v) \right\rangle}
      {|\g(u)-\g(v)|^2}}{|\g(u)-\g(v)|^2 }\right)\ |\g'(v)|\ |\g'(u)|\ dv\ du.
\end{align*} 
 
We will show that this formula is still valid under the weaker assumption that 
$\g \in H^{3/2}(\R / \Z,$ $\R^\tdim)$ with $\dg\in L^{\infty}$, i.e., for
arc-length parametrized $\g$ we only assume that the M\"obius energy
is finite.
We call a curve \emph{regular}
if there is a positive constant $c=c(\g)$ with
$|\g'(x)| \geq c$ {for all } $x \in \mathbb R/ \mathbb Z$.


\begin{theorem}[$E^{(2)}\in C^1(\Hir\capH)$] \label{thm:EulerLagrangeEquation}
 The energy $\E$ is continuously differentiable on the space of injective and regular
 curves belonging to $\W\capH$. Furthermore, if
 $\g\in\W$ is injective and parametrized by arc-length
 and $\h\in \W\capH$
%  \begin{enumerate}
%   \item there is some $\tau_0=\tau_0(\g,\varphi)>0$ such that
%    \[ \g+\tau\varphi\in\Wir \qquad\text{for all }\tau\in[-\tau_0,\tau_0], \]
%   \item
   the first variation \[ \delta\E(\g;\h) := \lim_{\tau\to0}\frac{\E(\g+\tau\h)-\E(\g)}\tau \]
   exists and equals
   \[%\begin{align}
     %\delta\E(\g;\h) &=
     2\lim_{\eps\searrow0} \iint\limits_{U_\eps}
     \br{\sp{\dg(u),\dh(u)}-\frac{\sp{\D\g,\D\h}}{\abs{\D\g}^2}} %\cdot \nonumber
     %&\qquad\qquad\qquad{}\cdot
     \frac{\d w\d u}{\abs{\D\g}^2}
   \]%\end{align}
   where
   \begin{equation}\label{eq:Ue}
     U_\eps := \R/\Z\times\br{[-\tfrac12,-\eps]\cup[\eps,\tfrac12]}.
   \end{equation}
% \end{enumerate}
\end{theorem}

Though the space $\W\capH$ seems somewhat artificial at first sight, it just guarantees that 
we do not use bad parametrizations of our curves. The proof of this 
result is similar to~\cite{BlattReiterStationary11}.

We will then use the resulting Euler-Langrange equation for stationary points of the M\"obius energy to
prove that these points are smooth:

\begin{theorem}[Stationary points are smooth] \label{thm:MainResult}
Any stationary point $\soln \in H^{\frac 3 2}(\R / \Z, \R^n)$ of $E^{(2)}$, i.e., any curve
$\soln \in H^{\frac 3 2 }(\R / \Z, \R^n)$ for which
\begin{equation*}
 \delta E^{(2)} (\g;h) =0 \quad \mbox{for all $h \in C^{\infty}(\mathbb R / \mathbb Z, \mathbb R^n)$},
\end{equation*}
belongs to $C^\infty$ when parametrized by arc-length.
\end{theorem}

In \cite{BlattReiterStationary11}, improving a previous result~\cite{ReiterOharaKnot11},
the smoothness of stationary points of finite energy was already shown
for the case of $E^{(\alpha)}:= E^{(\alpha,1)}$, $\alpha \in (2,3)$, instead of the M\"obius energy.
It is worth noting, that those energies lead to 
a \emph{subcritical} Euler-Langrage equation, and that in some sense the regularity theory can be based on  Sobolev 
embeddings for fractional Sobolev and Besov spaces. 
In contrast to this, the Euler-Lagrange equation of the M\"obius energy is \emph{critical}. As for 
well-known critical geometric equations -- like the Euler-Lagrange equation of the
Willmore functional, see, e.g., \cite{Simon1993, Riviere2006},
or harmonic maps on $\mathbb R^2$, see, e.g., \cite{Helein1991,Riviere2007}, -- one has first to find a way 
to gain an $\eps$ of additional regularity (via gaining a $\delta$ of additional integrability) and then start a bootstrapping argument.
That is, in a quite natural way, the proof of Theorem~\ref{thm:MainResult} is an immediate consequence of two technically independent steps:

\begin{theorem}[Initial regularity]\label{th:reg:1alpha} \label{thm:FirstStep}
Let $\soln \in H^{\frac{3}{2}}(\R /\Z,\R^\tdim)$, $\soln' \in \S^{\tdim-1}$, be a stationary point of the M\"obius energy, i.e. satisfying $\delta\E(\g,\varphi) = 0$ for all $\varphi\in C^\infty\rzd$.
Then $\soln \in C^{1,\alpha}$, for some $\alpha > 0$, and  $\soln \in H^{\frac{3}{2},p}(\R /\Z,\R^\tdim)$ for some $p > 2$.
\end{theorem}

\begin{theorem}[Bootstrapping] \label{thm:bootstrapping}
 For some $p > 2$, let $\soln \in H^{\frac{3}{2},p}(\R /\Z,\R^\tdim)$, $\soln' \in \S^{\tdim-1}$,
 be a stationary point of the M\"obius energy. Then $\soln$ is smooth.
\end{theorem}
Theorem~\ref{th:reg:1alpha} is proven in Section \ref{sec:initial}, Theorem~\ref{thm:bootstrapping} in Section \ref{sec:bootstrapping}.
While Theorem~\ref{thm:bootstrapping} relies mainly on bringing together Sobolev embeddings and 
standard commutator estimates for Bessel potential spaces with techniques developed in
\cite{Blatt2010a}, some very delicate estimates
are needed to get anything more than the critical and initial regularity $H^{\frac 3 2}$ for 
stationary points as stated in Theorem~\ref{thm:FirstStep}.

Both theorems rely on a decomposition of the first variation dating back to~\cite[Formula~(4.5)]{He2000} which 
already proved to be helpful in the analysis of the functionals 
$E^{(\alpha)}$ for $\alpha \in (2,3)$ (cf. \cite{BlattReiterStationary11}) and the gradient flow 
of the energies $E^{(\alpha)}$ for $\alpha \in [2,3)$ \cite{Blatt2011,Blatt2011c}. 

For  $f, g: \R \to \R^\tdim$, $\eps>0$ let 
\begin{equation}\label{eq:defQeps}
 Q_\varepsilon(f,g) := 
  \intl_{0}^1 \intl_{[-\frac{1}{2},\frac{1}{2}]\backslash (-\varepsilon,\varepsilon)} \br{\sbrac{{f}'(u), {g}'(u)}\ w^2 - \Sbrac{ f(u+w)-f(u), g(u+w)-g(u)}}
  \frac{dw}{w^4}\ du,
\end{equation}
\begin{equation}\label{eq:defT1}
 T_1(f,g) := - \intl_{0}^{1}\intl_{-\frac{1}{2}}^{\frac{1}{2}}\sbrac{{f}'(u),\ {g}'(u) }\ \brac{\fracm{\abs{f(u+w)-f(u)}^2}-\fracm{\abs{w}^2}}\ dw\ du,
\end{equation}
and
\begin{equation}\label{eq:defT2}
 T_2(f,g) := \intl_0^1\intl_{-\frac{1}{2}}^{\frac{1}{2}} \sbrac{ f(u+w)-f(u),\ g(u+w)-g(u) }\ \brac{\fracm{\abs{f(u+w)-f(u)}^4}-\fracm{\abs{w}^4}}\ dw\ du.
\end{equation}
From Theorem~\ref{thm:EulerLagrangeEquation} we deduce
that a critical knot $\soln \in H^{\frac{3}{2}}(\R /\Z,\R^\tdim)$,  parametrized by arc length, satisfies
\begin{equation}
\label{eq:knotPDE}
 Q(\soln,h) := \lim_{\varepsilon \to 0} Q_\varepsilon(\soln,h) = T_1(\soln,h) + T_2(\soln,h) \quad \mbox{for all $h \in C^\infty(\R /\Z)$}.
\end{equation}
This is the form of the Euler-Lagrange equation,  with which we will work in the proofs of Theorem~\ref{thm:FirstStep}
and Theorem~\ref{thm:bootstrapping}.

Let us conclude this introduction by remarking that, in contrast to stationary
points of $E^{(\alpha,1)}$, for $p>1$ we do not expect stationary points of $E^{(\alpha,p)}$ to be $C^\infty$-smooth:
The resulting Euler-Lagrange equation should be in some sense
a nonlocal degenerate elliptic equation. 
Keeping in mind the regularity theory for elliptic degenerate equations, one might expect nevertheless
that stationary points are at least 
a bit more regular than an arbitrary finite-energy curve alone. 



\paragraph*{Acknowledgements.}
The first author was supported by Swiss National Science Foundation Grant Nr.~200020\_125127 and
the Leverhulm trust.
The second author was supported by DFG Transregional Collaborative Research Centre SFB~TR~71.
The third author has received funding from the European Research Council under the European Union's Seventh Framework Programme (FP7/2007-2013) / ERC grant agreement no 267087, DAAD PostDoc Program (D/10/50763) and the Forschungsinstitut f\"ur Mathematik, ETH Z\"urich. He would like to thank Tristan Rivi\`ere and the ETH for their hospitality.